\section*{Erklärung zur Nutzung generativer Künstlicher Intelligenz}

Generative KI wurde im Verlauf dieser Masterarbeit regelmäßig als unterstützendes Werkzeug eingesetzt. Sie wurde als standardmäßiger Bestandteil moderner Arbeitsabläufe betrachtet -- in gleicher Weise, wie eine integrierte Entwicklungsumgebung oder eine Rechtschreibprüfung ein Standardbestandteil beim Schreiben von Programmcode und Fließtext ist. Der Einsatz erfolgte stets unterstützend und zu keinem Zeitpunkt als passive Abgabe von Aufgaben. Jede Ausgabe unterlag einer engen manuellen Überwachung, wurde wiederholt überarbeitet und durch den Autor verifiziert.

Die folgenden Werkzeuge kamen zum Einsatz:

\begin{itemize}
    \item \textbf{DeepL Write} wurde für das sprachliche Lektorat und das grammatikalische Korrekturlesen verwendet.
    \item \textbf{Gemini (Google)} wurde zur Unterstützung beim Brainstorming, zur Klärung technischer Konzepte und zur Generierung von Formulierungsvorschlägen für nicht-technische Erklärungstexte genutzt. Zudem diente es als erster Schritt bei der Informationssuche anstelle einer konventionellen Suchmaschine.
    \item \textbf{Antigravity (Google)} wurde als agentische Entwicklungsumgebung und über die Befehlszeilenschnittstelle (CLI) für Programmieraufgaben, technische Beratung, Code-Reviews sowie das Korrekturlesen der Abschlussarbeit eingesetzt (u.\,a. zur automatisierten Überprüfung auf argumentative Konsistenz, präzise Wortwahl sowie Rechtschreibung und Grammatik).
\end{itemize}

Die KI-Werkzeuge dienten zudem als Diskussions- und Sparringspartner sowie zum \enquote{Rubber Ducking}. Das strukturierte Formulieren eines Problems gegenüber dem System trug oft bereits vor Erhalt einer Antwort zur Klärung bei. Sämtliche methodischen Vorgehensweisen und erarbeiteten Lösungsansätze wurden darüber hinaus fortlaufend in den wöchentlichen Master-Meetings mit dem betreuenden Professor, den Doktoranden und weiteren Masterstudierenden gegengeprüft, diskutiert und entweder bestätigt, verworfen oder gezielt weiterentwickelt.

Alle durch KI-Werkzeuge generierten Inhalte wurden vom Autor kritisch geprüft, editiert und verifiziert. Kein generatives KI-System wurde zur Erzeugung originärer wissenschaftlicher Resultate, experimenteller Messdaten, Beweise oder der finalen Implementierungen der Algorithmen herangezogen. Das Forschungsdesign, die Analyse sowie die Interpretation der Ergebnisse sind originäre Leistungen des Autors. Der Autor übernimmt die uneingeschränkte Verantwortung für den Inhalt, die Richtigkeit und die Eigenständigkeit dieser Arbeit.
